\documentclass[12pt,a4paper]{article}

\usepackage[T1]{fontenc}
\usepackage[utf8]{inputenc}
\usepackage[brazil]{babel}
\usepackage{lmodern}
\usepackage{geometry}
\usepackage{hyperref}
\usepackage{enumitem}
\usepackage{array}

\geometry{margin=2.5cm}
\hypersetup{
  colorlinks=true,
  linkcolor=blue,
  urlcolor=blue
}

\title{Plano de Aula: Asserções e Gatilhos (Triggers)}
\date{}

\begin{document}
\maketitle

\noindent
\textbf{Duração:} 2 horas (120 minutos)\\
\textbf{Público-alvo:} Iniciantes\\
\textbf{Recursos Necessários:} Laboratório com computadores, SGBD e cliente SQL, projetor (opcional), quadro e marcadores.

\section{Objetivos da Aula}
\begin{itemize}[leftmargin=*,itemsep=0.25em]
  \item Compreender mecanismos de integridade e regras no banco de dados.
  \item Diferenciar restrições (PK/FK/UNIQUE/CHECK) de gatilhos (triggers).
  \item Implementar gatilhos simples para auditoria e validação.
  \item Reconhecer limitações e diferenças de suporte a asserções entre SGBDs.
\end{itemize}

\section{Metodologia}
Aula prática com exemplos progressivos. Primeiro, reforçar integridade com constraints. Depois, introduzir triggers como automações executadas em eventos (INSERT/UPDATE/DELETE). A atividade final aplica um trigger em um cenário realista.

\section{Cronograma e Conteúdo Detalhado (120 Minutos)}
\subsection*{Parte 1: Integridade e Regras (20 minutos)}
\begin{itemize}[leftmargin=*,itemsep=0.35em]
  \item Revisar constraints: PK, FK, UNIQUE e CHECK.
  \item Analogia: ``catracas'' que impedem entrar dado inválido.
\end{itemize}

\subsection*{Parte 2: Asserções (15 minutos)}
\begin{itemize}[leftmargin=*,itemsep=0.35em]
  \item Conceito: regra global que deve ser verdadeira para o banco.
  \item Discutir suporte do SGBD adotado e alternativas (VIEW + CHECK/trigger).
\end{itemize}

\subsection*{Parte 3: Triggers --- Conceito e Estrutura (25 minutos)}
\begin{itemize}[leftmargin=*,itemsep=0.35em]
  \item O que é trigger: ``uma regra automática que dispara em um evento''.
  \item Eventos: BEFORE/AFTER INSERT/UPDATE/DELETE (conforme SGBD).
  \item Usos comuns: auditoria, validação complexa, manutenção de históricos.
\end{itemize}

\subsection*{Parte 4: Exemplo Guiado (25 minutos)}
\begin{itemize}[leftmargin=*,itemsep=0.35em]
  \item Criar tabela de auditoria (ex.: LOG\_PET).
  \item Criar trigger que registra alterações de PET (UPDATE/DELETE).
  \item Testar e interpretar o log gerado.
\end{itemize}

\subsection*{Parte 5: Atividade em Duplas --- Regra Automática (30 minutos)}
\begin{itemize}[leftmargin=*,itemsep=0.35em]
  \item Criar um trigger para:
  \begin{itemize}[leftmargin=*,itemsep=0.25em]
    \item Impedir atualização inválida (ex.: campo vazio).
    \item Ou manter histórico de alterações em uma tabela de log.
  \end{itemize}
  \item Entrega: script do trigger + evidência de execução (antes/depois).
\end{itemize}

\subsection*{Parte 6: Encerramento (5 minutos)}
\begin{itemize}[leftmargin=*,itemsep=0.35em]
  \item Revisão: quando usar constraint vs. trigger.
  \item Ponte: ``Próxima aula: controle de transações e consistência (ACID).''
\end{itemize}

\end{document}

